\documentclass{beamer}

\title{Mobile Scanner Application}
\author{Bhoomika L -EE25BTECH11014\\Vaishnavi R.A-EE25BTECH11059}
\date{}

\begin{document}

\frame{\titlepage}

%--------------------------------
\begin{frame}{What does the system do?}
\begin{itemize}
\item This system converts a mobile camera into a document scanner.
\item It automatically detects a paper in the camera view.
\item The document is corrected and saved like a scanned copy.
\item Multiple pages can be combined into a PDF.
\end{itemize}
\end{frame}

%--------------------------------
\begin{frame}{Modes of Operation}
\begin{itemize}
\item The system works in two modes:
\item Continuous Scan Mode (Video Mode)
\item Multi Document Scan Mode
\item Both modes use the same detection pipeline.
\item The difference lies in how capturing is triggered.
\end{itemize}
\end{frame}

%--------------------------------
\begin{frame}{Continuous Scan Mode (Idea)}
\begin{itemize}
\item Detects one document at a time.
\item Once the document is stable and clear, it is captured.
\item Useful for scanning pages one by one.
\end{itemize}
\end{frame}

%--------------------------------
\begin{frame}{Continuous Mode - Logic Flow}
\begin{itemize}
\item Capture frame from camera
\item Detect rectangle in frame
\item Apply perspective correction
\item Check focus using variance
\item Start timer when stable
\item Capture image after fixed time
\end{itemize}
\end{frame}

%--------------------------------
\begin{frame}{Multi Document Mode (Idea)}
\begin{itemize}
\item Detects multiple documents at the same time.
\item Observes the scene for a few seconds.
\item Captures all documents together when stable.
\end{itemize}
\end{frame}

%--------------------------------
\begin{frame}{Multi Mode - Logic Flow}
\begin{itemize}
\item Capture frame from camera
\item Detect all rectangles
\item Count number of valid documents
\item Track maximum count
\item Wait until count becomes stable
\item Capture all documents
\end{itemize}
\end{frame}

%--------------------------------
\begin{frame}{Overall Processing Pipeline}
\begin{itemize}
\item Frame Capture → Preprocessing → Edge Detection
\item Contour Detection → Document Selection
\item Perspective Transformation → Focus Check → Capture
\end{itemize}
\end{frame}

%--------------------------------
\begin{frame}{Why Preprocessing?}
\begin{itemize}
\item Raw images contain noise and unnecessary color information.
\item These make detection difficult.
\item So we simplify the image before processing.
\end{itemize}
\end{frame}

%--------------------------------
\begin{frame}{Grayscale Conversion}
\begin{itemize}
\item Converts RGB image to grayscale.
\item Reduces complexity from 3 channels to 1.
\item Edge detection depends on intensity, not color.
\end{itemize}
\end{frame}

%--------------------------------
\begin{frame}{Noise Removal}
\begin{itemize}
\item Noise affects detection accuracy.
\item Bilateral filtering removes noise.
\item Edges are preserved for better detection.
\end{itemize}
\end{frame}

%--------------------------------
\begin{frame}{Edge Detection}
\begin{itemize}
\item Edges are regions where pixel intensity changes sharply.
\item These correspond to boundaries of objects.

\item Canny edge detection is used:
\begin{itemize}
\item Computes intensity gradients
\item Uses lower and upper thresholds
\end{itemize}

\item Decision:
\begin{itemize}
\item Above upper → strong edge
\item Between → weak edge
\item Below → ignored
\end{itemize}
\end{itemize}
\end{frame}

%--------------------------------
\begin{frame}{Adaptive Thresholding}
\begin{itemize}
\item Lighting conditions may vary.
\item Thresholds are computed using median intensity.
\end{itemize}

\[
v = \text{median intensity}
\]

\begin{itemize}
\item Lower = $0.66v$
\item Upper = $1.33v$
\end{itemize}
\end{frame}

%--------------------------------
\begin{frame}{Morphological Processing}
\begin{itemize}
\item Dilation connects broken edges.
\item Closing fills gaps.
\item Ensures continuous boundaries.
\end{itemize}
\end{frame}

%--------------------------------
\begin{frame}{Contour Detection}
\begin{itemize}
\item Contours represent object boundaries.
\item All shapes in the image are detected.
\end{itemize}
\end{frame}

%--------------------------------
\begin{frame}{Filtering and Shape Detection}
\begin{itemize}
\item Small contours are removed.
\item Shapes with 4 corners are selected.
\end{itemize}
\end{frame}

%--------------------------------
\begin{frame}{Selecting the Document}
\begin{itemize}
\item Largest rectangle is selected.
\item Assumed to be the document.
\end{itemize}
\end{frame}

%--------------------------------
\begin{frame}{Perspective Transformation}
\begin{itemize}
\item Converts tilted document into flat image.
\end{itemize}

\[
\begin{bmatrix}x'\\y'\\1\end{bmatrix}
=
H
\begin{bmatrix}x\\y\\1\end{bmatrix}
\]
\end{frame}

%--------------------------------
\begin{frame}{Focus Detection}
\begin{itemize}
\item Uses variance of Laplacian.
\item High variance → sharp
\item Low variance → blur
\item Threshold = 120
\end{itemize}
\end{frame}

%--------------------------------
\begin{frame}{Auto Capture}
\begin{itemize}
\item Captures when document is stable.
\item Uses time-based triggering.
\end{itemize}
\end{frame}
\begin{frame}{Duplicate Detection}
\begin{itemize}
\item Compares current and previous scan.
\item Uses image difference:
\[
\text{score} = \text{mean difference}
\]
\item Low score → duplicate → skipped
\end{itemize}
\end{frame}

%--------------------------------
\begin{frame}{PDF Creation}
\begin{itemize}
\item Images are combined into PDF.
\end{itemize}
\end{frame}

%================= EXTRA FEATURES =================





\begin{frame}{Observation Phase}
\begin{itemize}
\item System observes scene for few seconds.
\item Tracks maximum document count.
\item Ensures all documents are detected.
\end{itemize}
\end{frame}

\begin{frame}{Stability Check}
\begin{itemize}
\item Capture only when document is stable.
\item Reduces blur and improves quality.
\end{itemize}
\end{frame}

\begin{frame}{Efficient Comparison}
\begin{itemize}
\item Images resized before comparison.
\item Faster and efficient processing.
\end{itemize}
\end{frame}


%--------------------------------


\end{document}